\documentclass{situationist-preprint}

\title{Impossible Worlds and Hyperintensionality}
\msid{TSM-09-a}
\msfamily{TSM-09}
\msdiscipline{Time, Space and Modality}
\mstype{Principal manuscript}
\mskeywords{impossible worlds; hyperintensionality; content; logical omniscience; counterpossibles; ersatzism}

\begin{document}
\maketitle

\begin{abstract}
Are there impossible worlds? We argue that there are, because they are required to model hyperintensional content: contents that differ although they are true in exactly the same possible worlds. Belief, information and counterpossible supposition discriminate between them. We prove that any model on which contents are sets of points of evaluation, and on which necessarily equivalent contents can differ, must contain points at which some necessary truth fails or some impossibility holds; such points are impossible worlds. The point format is not optional: it gives content its logic of entailment and update, which structured propositions do not supply. Impossible worlds, constructed as in TSM-07 from propositions, are sets that are not consistent. They are not possibilities: the modal quantifiers range over the possible worlds only. They are as innocent as possible worlds and confined to the theory of content.
\end{abstract}

\section{The problem}

On the intensional theory of content, developed by \citet{Stalnaker1984}, the content of a sentence or a belief is the set of possible worlds at which it is true. The theory has a well-known consequence: necessarily equivalent contents are identical. Every necessary truth has the same content, the set of all worlds, and so does every mathematical truth. To believe that $2+2=4$ is to believe that there are infinitely many primes; a subject who believes the premises of a valid argument believes its conclusion. This is the problem of logical omniscience \citep{Hintikka1975}. Stalnaker located the failures of omniscience in metalinguistic ignorance, which is unconvincing for mathematical belief.

The phenomenon is wider than belief. A context is \emph{hyperintensional} if substitution of necessarily equivalent expressions in it can change truth value \citep{Cresswell1975}. Attitude reports are hyperintensional; so are explanations; so are counterpossible conditionals, of which some are true and others false though their antecedents are equally impossible \citep{Nolan1997}.

Two responses are in the literature. One introduces structured propositions, which are individuated by their constituents and the way they are combined \citep{Soames1987}. The other enlarges the space of points with impossible worlds, at which necessary truths may fail \citep{Nolan1997,Jago2014,BertoJago2019}. We argue that the second is required.

\section{Framework}

\begin{definition}[Point model of content]\label{d:model}
A \emph{point model} of content is a set $P$ of points of evaluation, including the possible worlds, together with an assignment to each content $c$ of the set $|c|\subseteq P$ of points at which it holds. Contents $c$ and $d$ are \emph{distinguished} by the model if $|c|\ne|d|$.
\end{definition}

\begin{definition}[Necessary equivalence]\label{d:equiv}
Contents are \emph{necessarily equivalent} if they hold at the same possible worlds.
\end{definition}

\begin{definition}[Impossible world]\label{d:imp}
An \emph{impossible world} is a point of evaluation at which the contents that hold are not those that hold at any possible world.
\end{definition}

On the account of TSM-07, possible worlds are maximal consistent sets of propositions, a proposition holding at a world when it is a member. Impossible worlds are sets of propositions that are not consistent, whether maximal or not: sets that contain a proposition and its negation, or that contain propositions logically consistent with each other but inconsistent with the essences of things, in the sense of TSM-10. Possibility is defined, in TSM-07, by quantification over possible worlds only; impossible worlds therefore do not make anything possible.

\begin{principle}[Hyperintensionality]\label{p:hyper}
Some necessarily equivalent contents are distinct, and an adequate model of content distinguishes them.
\end{principle}

\begin{principle}[Point format]\label{p:format}
An adequate model of content is a point model.
\end{principle}

\section{The argument}

\begin{theorem}\label{t:main}
Every adequate model of content contains impossible worlds.
\end{theorem}

\begin{proof}
By Principles~\ref{p:hyper} and~\ref{p:format}, an adequate model is a point model that distinguishes some necessarily equivalent contents $c$ and $d$. Since $|c|\ne|d|$, there is a point $x$ at which one holds and the other does not, say $c$ but not $d$. Since $c$ and $d$ hold at the same possible worlds, at every possible world either both hold or neither does; at $x$, one holds and the other does not. So the contents holding at $x$ are not those holding at any possible world, and $x$ is an impossible world.
\end{proof}

The work is in the principles. Principle~\ref{p:hyper} is the datum of Section~1; Principle~\ref{p:format} is what the structured-proposition theorist rejects.

\section{Why the point format}

Contents do more than differ. They stand in relations of entailment, they combine by conjunction and disjunction, and they are added to bodies of information and subtracted from them. In a point model, these relations are set-theoretic: $c$ entails $d$ when $|c|\subseteq|d|$, conjunction is intersection, and update with $c$ is intersection with $|c|$. The theories of conversation and belief revision use these operations \citep{Stalnaker1984}.

Structured propositions individuate content, but they supply no such operations: that the structured propositions that $2+2=4$ and that there are infinitely many primes differ does not tell us what a subject who accepts one and not the other takes the world to be like. To say this, she must say which situations each structured proposition rules out; she then has a point model, and by Theorem~\ref{t:main} the situations include impossible ones \citep{Jago2014,BertoJago2019}.

The same holds for counterpossibles: on the similarity semantics, a conditional with an impossible antecedent is vacuously true unless there are impossible points at which the antecedent holds, ordered by similarity to the actual world \citep{Nolan1997}.

\section{Objections}

\begin{objection}[Impossible worlds are impossible]
An impossible world is a world at which impossibilities obtain. To accept that there are such worlds is to accept that impossibilities obtain somewhere, which is to accept impossibilities.
\end{objection}

\begin{reply}
It succeeds only against concrete impossible worlds. There are no concrete worlds other than the actual one (TSM-08); impossible worlds, like possible ones, are sets of propositions. An impossible world at which something is round and square is a set containing that proposition; it is not round or square, as MET-24 also observes.
\end{reply}

\begin{objection}[Modal inflation]
If there are impossible worlds, there are worlds at which a contradiction is true. Then a contradiction is true at some world, and so possible.
\end{objection}

\begin{reply}
Possibility is truth at some possible world (TSM-07), and necessity truth at every world consistent with the essences (TSM-10). Impossible worlds are not in the domain of the modal quantifiers, and nothing is made possible by holding at them.
\end{reply}

\begin{objection}[Anything goes]
If impossible worlds are arbitrary sets of propositions, content becomes as fine as the propositions themselves, and the point format adds nothing.
\end{objection}

\begin{reply}
The point format adds the operations, which individuation alone does not; and which impossible worlds a model uses depends on its purpose, the space of all sets of propositions being the domain from which models select. The difficulty \citet{Bjerring2013} raised for models of non-ideal reasoners is a problem of selection, discussed in TSM-09-b, not of the existence of the points.
\end{reply}

\section{Conclusion}

Hyperintensional content requires a point model in which necessarily equivalent contents differ, and any such model contains impossible worlds: sets of propositions that are not consistent, confined to the theory of content, the modal quantifiers ranging over the possible worlds alone.

\begin{thebibliography}{99}\small
\bibitem[Berto and Jago(2019)]{BertoJago2019} Berto, F., and Jago, M. (2019). \emph{Impossible Worlds}. Oxford University Press.
\bibitem[Bjerring(2013)]{Bjerring2013} Bjerring, J.~C. (2013). Impossible worlds and logical omniscience: An impossibility result. \emph{Synthese}, 190, 2505--2524.
\bibitem[Cresswell(1975)]{Cresswell1975} Cresswell, M.~J. (1975). Hyperintensional logic. \emph{Studia Logica}, 34, 25--38.
\bibitem[Hintikka(1975)]{Hintikka1975} Hintikka, J. (1975). Impossible possible worlds vindicated. \emph{Journal of Philosophical Logic}, 4, 475--484.
\bibitem[Jago(2014)]{Jago2014} Jago, M. (2014). \emph{The Impossible: An Essay on Hyperintensionality}. Oxford University Press.
\bibitem[Nolan(1997)]{Nolan1997} Nolan, D. (1997). Impossible worlds: A modest approach. \emph{Notre Dame Journal of Formal Logic}, 38(4), 535--572.
\bibitem[Soames(1987)]{Soames1987} Soames, S. (1987). Direct reference, propositional attitudes, and semantic content. \emph{Philosophical Topics}, 15, 47--87.
\bibitem[Stalnaker(1984)]{Stalnaker1984} Stalnaker, R. (1984). \emph{Inquiry}. MIT Press.
\end{thebibliography}

\end{document}
